%%
%% This is file `t-markdownthemewitiko_markdown_defaults.tex',
%% generated with the docstrip utility.
%%
%% The original source files were:
%%
%% markdown.dtx  (with options: `themes-witiko-markdown-defaults-ctx')
%% 
%% Copyright (C) 2016-2025 Vít Starý Novotný, Andrej Genčur
%% 
%% This work may be distributed and/or modified under the
%% conditions of the LaTeX Project Public License, either version 1.3c
%% of this license or (at your option) any later version.
%% The latest version of this license is in
%% 
%%    http://www.latex-project.org/lppl.txt
%% 
%% This work has the LPPL maintenance status `maintained'.
%% The Current Maintainer of this work is Vít Starý Novotný.
%% 
%% Send bug reports, requests for additions and questions
%% either to the GitHub issue tracker at
%% 
%%   https://github.com/Witiko/markdown/issues
%% 
%% or to the e-mail address <witiko@mail.muni.cz>.
%% 
%% MODIFICATION ADVICE:
%% 
%% If you want to customize this file, it is best to make a copy of
%% the source file(s) from which it was produced.  Use a different
%% name for your copy(ies) and modify the copy(ies); this will ensure
%% that your modifications do not get overwritten when you install a
%% new release of the standard system.  You should also ensure that
%% your modified source file does not generate any modified file with
%% the same name as a standard file.
%% 
%% You will also need to produce your own, suitably named, .ins file to
%% control the generation of files from your source file; this file
%% should contain your own preambles for the files it generates, not
%% those in the standard .ins files.
%% 
%% The names of the source files used are shown above.
%% 
\startmodule[markdownthemewitiko_markdown_defaults]
\unprotect
\markdownLoadPlainTeXTheme
\markdownIfOption{plain}{\iffalse}{\iftrue}
\def\markdownRendererHardLineBreakPrototype{\blank}%
\def\markdownRendererLeftBracePrototype{\textbraceleft}%
\def\markdownRendererRightBracePrototype{\textbraceright}%
\def\markdownRendererDollarSignPrototype{\textdollar}%
\def\markdownRendererPercentSignPrototype{\percent}%
\def\markdownRendererUnderscorePrototype{\textunderscore}%
\def\markdownRendererCircumflexPrototype{\textcircumflex}%
\def\markdownRendererBackslashPrototype{\textbackslash}%
\def\markdownRendererTildePrototype{\textasciitilde}%
\def\markdownRendererPipePrototype{\char`|}%
\def\markdownRendererLinkPrototype#1#2#3#4{%
  \useURL[#1][#3][][#4]#1\footnote[#1]{\ifx\empty#4\empty\else#4:
  \fi\tt<\hyphenatedurl{#3}>}}%
\usemodule[database]
\defineseparatedlist
  [MarkdownConTeXtCSV]
  [separator={,},
   before=\bTABLE,after=\eTABLE,
   first=\bTR,last=\eTR,
   left=\bTD,right=\eTD]
\def\markdownConTeXtCSV{csv}
\def\markdownRendererContentBlockPrototype#1#2#3#4{%
  \def\markdownConTeXtCSV@arg{#1}%
  \ifx\markdownConTeXtCSV@arg\markdownConTeXtCSV
    \placetable[][tab:#1]{#4}{%
      \processseparatedfile[MarkdownConTeXtCSV][#3]}%
  \else
    \markdownInput{#3}%
  \fi}%
\def\markdownRendererImagePrototype#1#2#3#4{%
  \placefigure[][]{#4}{\externalfigure[#3]}}%
\def\markdownRendererUlBeginPrototype{\startitemize}%
\def\markdownRendererUlBeginTightPrototype{\startitemize[packed]}%
\def\markdownRendererUlItemPrototype{\item}%
\def\markdownRendererUlEndPrototype{\stopitemize}%
\def\markdownRendererUlEndTightPrototype{\stopitemize}%
\def\markdownRendererOlBeginPrototype{\startitemize[n]}%
\def\markdownRendererOlBeginTightPrototype{\startitemize[packed,n]}%
\def\markdownRendererOlItemPrototype{\item}%
\def\markdownRendererOlItemWithNumberPrototype#1{\sym{#1.}}%
\def\markdownRendererOlEndPrototype{\stopitemize}%
\def\markdownRendererOlEndTightPrototype{\stopitemize}%
\definedescription
  [MarkdownConTeXtDlItemPrototype]
  [location=hanging,
   margin=standard,
   headstyle=bold]%
\definestartstop
  [MarkdownConTeXtDlPrototype]
  [before=\blank,
   after=\blank]%
\definestartstop
  [MarkdownConTeXtDlTightPrototype]
  [before=\blank\startpacked,
   after=\stoppacked\blank]%
\def\markdownRendererDlBeginPrototype{%
  \startMarkdownConTeXtDlPrototype}%
\def\markdownRendererDlBeginTightPrototype{%
  \startMarkdownConTeXtDlTightPrototype}%
\def\markdownRendererDlItemPrototype#1{%
  \startMarkdownConTeXtDlItemPrototype{#1}}%
\def\markdownRendererDlItemEndPrototype{%
  \stopMarkdownConTeXtDlItemPrototype}%
\def\markdownRendererDlEndPrototype{%
  \stopMarkdownConTeXtDlPrototype}%
\def\markdownRendererDlEndTightPrototype{%
  \stopMarkdownConTeXtDlTightPrototype}%
\def\markdownRendererEmphasisPrototype#1{{\em#1}}%
\def\markdownRendererStrongEmphasisPrototype#1{{\bf#1}}%
\def\markdownRendererBlockQuoteBeginPrototype{\startquotation}%
\def\markdownRendererBlockQuoteEndPrototype{\stopquotation}%
\def\markdownRendererLineBlockBeginPrototype{%
  \begingroup
    \def\markdownRendererHardLineBreak{
    }%
    \startlines
}%
\def\markdownRendererLineBlockEndPrototype{%
    \stoplines
  \endgroup
}%
\def\markdownRendererInputVerbatimPrototype#1{\typefile{#1}}%
\ExplSyntaxOn
\cs_gset:Npn
  \markdownRendererInputFencedCodePrototype#1#2#3
  {
    \tl_if_empty:nTF
      { #2 }
      { \markdownRendererInputVerbatim{#1} }
      {
        \regex_extract_once:nnN
          { \w* }
          { #2 }
          \l_tmpa_seq
        \seq_pop_left:NN
          \l_tmpa_seq
          \l_tmpa_tl
        \typefile[ \l_tmpa_tl ][] {#1}
      }
  }
\ExplSyntaxOff
\def\markdownRendererHeadingOnePrototype#1{\chapter{#1}}%
\def\markdownRendererHeadingTwoPrototype#1{\section{#1}}%
\def\markdownRendererHeadingThreePrototype#1{\subsection{#1}}%
\def\markdownRendererHeadingFourPrototype#1{\subsubsection{#1}}%
\def\markdownRendererHeadingFivePrototype#1{\subsubsubsection{#1}}%
\def\markdownRendererHeadingSixPrototype#1{\subsubsubsubsection{#1}}%
\def\markdownRendererThematicBreakPrototype{%
  \blackrule[height=1pt, width=\hsize]}%
\def\markdownRendererNotePrototype#1{\footnote{#1}}%
\def\markdownRendererTickedBoxPrototype{$\boxtimes$}
\def\markdownRendererHalfTickedBoxPrototype{$\boxdot$}
\def\markdownRendererUntickedBoxPrototype{$\square$}
\def\markdownRendererStrikeThroughPrototype#1{\overstrikes{#1}}
\def\markdownRendererSuperscriptPrototype#1{\high{#1}}
\def\markdownRendererSubscriptPrototype#1{\low{#1}}
\def\markdownRendererDisplayMathPrototype#1{%
  \startformula#1\stopformula}%
\newcount\markdownConTeXtRowCounter
\newcount\markdownConTeXtRowTotal
\newcount\markdownConTeXtColumnCounter
\newcount\markdownConTeXtColumnTotal
\newtoks\markdownConTeXtTable
\newtoks\markdownConTeXtTableFloat
\def\markdownRendererTablePrototype#1#2#3{%
  \markdownConTeXtTable={}%
  \ifx\empty#1\empty
    \markdownConTeXtTableFloat={%
      \the\markdownConTeXtTable}%
  \else
    \markdownConTeXtTableFloat={%
      \placetable{#1}{\the\markdownConTeXtTable}}%
  \fi
  \begingroup
  \setupTABLE[r][each][topframe=off, bottomframe=off,
                       leftframe=off, rightframe=off]
  \setupTABLE[c][each][topframe=off, bottomframe=off,
                       leftframe=off, rightframe=off]
  \setupTABLE[r][1][topframe=on, bottomframe=on]
  \setupTABLE[r][#1][bottomframe=on]
  \markdownConTeXtRowCounter=0%
  \markdownConTeXtRowTotal=#2%
  \markdownConTeXtColumnTotal=#3%
  \markdownConTeXtRenderTableRow}
\def\markdownConTeXtRenderTableRow#1{%
  \markdownConTeXtColumnCounter=0%
  \ifnum\markdownConTeXtRowCounter=0\relax
    \markdownConTeXtReadAlignments#1%
    \markdownConTeXtTable={\bTABLE}%
  \else
    \markdownConTeXtTable=\expandafter{%
      \the\markdownConTeXtTable\bTR}%
    \markdownConTeXtRenderTableCell#1%
    \markdownConTeXtTable=\expandafter{%
      \the\markdownConTeXtTable\eTR}%
  \fi
  \advance\markdownConTeXtRowCounter by 1\relax
  \ifnum\markdownConTeXtRowCounter>\markdownConTeXtRowTotal\relax
    \markdownConTeXtTable=\expandafter{%
      \the\markdownConTeXtTable\eTABLE}%
    \the\markdownConTeXtTableFloat
    \endgroup
    \expandafter\gobbleoneargument
  \fi\markdownConTeXtRenderTableRow}
\def\markdownConTeXtReadAlignments#1{%
  \advance\markdownConTeXtColumnCounter by 1\relax
  \if#1d%
    \setupTABLE[c][\the\markdownConTeXtColumnCounter][align=right]
  \fi\if#1l%
    \setupTABLE[c][\the\markdownConTeXtColumnCounter][align=right]
  \fi\if#1c%
    \setupTABLE[c][\the\markdownConTeXtColumnCounter][align=middle]
  \fi\if#1r%
    \setupTABLE[c][\the\markdownConTeXtColumnCounter][align=left]
  \fi
  \ifnum\markdownConTeXtColumnCounter<\markdownConTeXtColumnTotal\relax
  \else
    \expandafter\gobbleoneargument
  \fi\markdownConTeXtReadAlignments}
\def\markdownConTeXtRenderTableCell#1{%
  \advance\markdownConTeXtColumnCounter by 1\relax
  \markdownConTeXtTable=\expandafter{%
    \the\markdownConTeXtTable\bTD#1\eTD}%
  \ifnum\markdownConTeXtColumnCounter<\markdownConTeXtColumnTotal\relax
  \else
    \expandafter\gobbleoneargument
  \fi\markdownConTeXtRenderTableCell}
\ExplSyntaxOn
\cs_gset:Npn
  \markdownRendererInputRawInlinePrototype#1#2
  {
    \str_case:nnF
      { #2 }
      {
        { latex }
          {
            \__markdown_plain_tex_default_input_raw_inline:nn
              { #1 }
              { context }
          }
      }
      {
        \__markdown_plain_tex_default_input_raw_inline:nn
          { #1 }
          { #2 }
      }
  }
\cs_gset:Npn
  \markdownRendererInputRawBlockPrototype#1#2
  {
    \str_case:nnF
      { #2 }
      {
        { context }
          {
            \__markdown_plain_tex_default_input_raw_block:nn
              { #1 }
              { tex }
          }
      }
      {
        \__markdown_plain_tex_default_input_raw_block:nn
          { #1 }
          { #2 }
      }
  }
\cs_gset_eq:NN
  \markdownRendererInputRawBlockPrototype
  \markdownRendererInputRawInlinePrototype
\fi % Closes `\markdownIfOption{plain}{\iffalse}{\iftrue}`
\ExplSyntaxOff
\stopmodule
\protect
\endinput
%%
%% End of file `t-markdownthemewitiko_markdown_defaults.tex'.
